\section{Weighted models and the reduction from arbitrary Kakeya sets}
\label{sec:framework}

We encode a line family by a probability law on trajectories and a
selection of trajectory--time incidences. A local fit confines the
base spatial projection of each whole trajectory piece to a moving
box and keeps a quadratic polynomial small along that piece. Its mass is measured on the
selected incidences.
Theorem~\ref{thm:model-reduction} turns a hypothetical Kakeya
counterexample into a weighted problem in which fitting at the final
thickness requires a positive power cost in time localization.
The infimum of these costs is the horizon.

We first develop the observations and density estimates used to measure
these fits. We then define the charts, show how their priors and densities
transform under normalization, and prove the reduction from an arbitrary
Kakeya set. All subpower factors in this section refer to one exact
sequence with resolution tending to infinity; later sequences of exact
problems will be introduced separately.

\subsection{Exact problems and observations}

The model has two versions. Version \(2\) retains the lines in
\(\R^4\); version \(1\) is the scalar model used after projection.
In each version a prior samples a trajectory independently of time,
and a separate event specifies which trajectory--time incidences
are retained.

\begin{definition}[Exact problems]\label{def:exact-problem}
An exact problem is a sequence indexed by a real parameter
\(N_0\to\infty\), with a fixed length \(0<L<\infty\).
For \(0\le r\le L\), write \(a_r=N_0^{-r}\).
A subpower factor \(K_0\ge1\) satisfies
\(\log K_0/\log N_0\to0\); its value may be increased from one occurrence
to the next.  We write \(A\sim_{\log}B\) when
\(K_0^{-1}B\le A\le K_0B\) for such a factor.
Subsequences are permitted.

There are finitely many auxiliary worlds \(\gamma\), of weights
\(\omega_\gamma\ge0\) satisfying \(\sum_\gamma\omega_\gamma=1\).
In world \(\gamma\), a trajectory has prior probability law
\(\nu_\gamma\), independently of uniform time \(t\in[0,1]\).
Its base coordinate is
\[
 y(t)=y_0+tV\in\R^j,\qquad |y(t)|+|V|\le K_0.
\]
We use the following two versions.
\begin{enumerate}
\item In version \(1\), \(j=1\), and
\(w(t)=w_0+w_1t+w_2t^2\), with
\(\sup_{[0,1]}|w|\le K_0\).
Tests are \(P(t,y,w)=w-F(t,y)\), where \(F\) is a polynomial of
total degree at most two.  The derivative scale is \(B=1\), and \(X=w\).
\item In version \(2\), \(j=2\), and \(w(t)=w_0+tW\) is affine.
The coefficients of \(w\) need not be independent of those of \(y\).
Each world has a derivative scale \(B>0\) and a polynomial
\(P_*(t,y,w)\) of total degree at most two.  Writing
\(X(t)=P_*(t,y(t),w(t))\), the native bounds are
\begin{align}
 \sup_{[0,1]}|X|&\le K_0,&
 \sup_{[0,1]}|\partial_wP_*|&\le K_0B,
 \label{a01:native-bounds}\\
 |\partial_wP_*|&\ge B/K_0
       \quad\hbox{at selected incidences},&
 |\partial_{ww}P_*|&\le K_0B^2.
 \notag
\end{align}
Tests in this version are arbitrary polynomials \(P(t,y,w)\) of
total degree at most two.
\end{enumerate}
In both versions \(X\) is a quadratic polynomial along each trajectory;
its coefficients are subpower bounded by interpolation on \([0,1]\).

An incidence selection is an event \(E_\gamma\) in trajectory--time
space, and
\begin{equation}
 d\mu_\gamma(l,t)=\1_{E_\gamma}(l,t)\,d\nu_\gamma(l)\,dt,
 \qquad
 \mu=\sum_\gamma\omega_\gamma\mu_\gamma,\qquad
 \mu(\mathrm{all})\ge K_0^{-1}.
 \label{a01:incidence-law}
\end{equation}
Thus \(\mu_\gamma\) is not normalized to probability.
A further selection restricts this event and must retain total mass
at least an inverse subpower.
\end{definition}

Here and below polynomial bounds are \(N_0^M\) for a finite exponent
\(M\) fixed throughout the exact sequence.  The exponent may be
arbitrarily large and may change when a new exact problem is constructed.
We impose the following finiteness convention.

\begin{definition}[Tempered data]\label{a01:tempered}
The numbers of worlds, specified parameter ranges and coefficients,
and \(B,B^{-1}\), are polynomially bounded.  In the full trajectory
coefficient coordinates, each prior has piecewise constant Lebesgue
density, polynomially bounded on polynomially many semialgebraic
pieces.  Trajectory and incidence predicates use polynomially many
polynomial conditions of bounded individual degree, with arbitrary
Boolean combinations.  The number of variables in each geometric
test is bounded.  Real constants in these conditions have no height
restriction.

Exact-level selections, labels, cells, and frames obey such uniform
bounds.  Chart widths and time lengths have polynomially bounded
reciprocals.  Changes of trajectory coordinates are invertible and
have polynomially controlled constant Jacobian.  We permit
polynomially many copies of coefficient space, with an internal
discrete trajectory index.  This index is not a world or an observed
point feature; all assignments concern indexed trajectories.

Measurable arguments may be used to locate a finite family of boxes
or polynomial tests.  The output is encoded by those constants and
geometric tests, not by a description of the intermediate measurable
search.
\end{definition}

Worlds with extremely small success fractions can be removed.
After forming chart worlds, charts with extremely small conditional
prior mass can also be removed: there are only polynomially many
charts, and the selected measure is bounded by prior times time.
The threshold can be a sufficiently small reciprocal polynomial.
These observations preserve the convention in
\Cref{a01:tempered}.

Use nested half-open dyadic grids, rounding requested widths upwards
within a factor of two.  The observation at depth \(r\) is
\begin{equation}
 p_r=(\gamma,t,y(t),[w(t)]_{a_r/B}).
 \label{a01:observation}
\end{equation}
Within a world its density is with respect to \(dt\,dy\) and counting
measure on the hidden bin.  Such densities exist: at fixed \(t\),
replacing \(y_0\) by \(y(t)=y_0+tV\) has Jacobian one.

A density is \emph{typically} \(n=N_0^{-f}\) if, for every fixed
\(\tau>0\), the selected incidences where it is outside
\([N_0^{-\tau}n,N_0^\tau n]\) have total weighted mass at most
\(N_0^{-c(\tau)}\), eventually, for some \(c(\tau)>0\).
One-sided typical bounds have the corresponding meaning.
A problem is homogeneous for a family of observations if each
observation has a common exponent, independent of the world and
fixed along the exact sequence.  Unless otherwise stated the
observations are \(p_r\), \(0\le r\le L\), and their typical densities
are denoted \(n(r)=N_0^{-f(r)}\).

\subsection{Density selection and homogeneous profiles}

The selected observation densities need not initially have uniform
sizes. We will restrict to incidences on which their logarithmic sizes
stabilize. Restrictions and finite refinements obey a simple mass
estimate; a finite-table approximation will make the subsequent density
selections compatible with the tempered data.

\begin{lemma}[Restriction, refinement, and caps]
\label{a01:density-operations}
Let \(g\) be the unnormalized density of an observation and \(g'\)
the density after restriction.  For \(0<\epsilon<1\),
\begin{equation}
 \int_{\{g'<\epsilon g\}}g'\le
 \epsilon\int g.
 \label{a01:restriction-integral}
\end{equation}
If an observation is refined by a discrete label with at most \(D\)
possible values on the retained events, the total refined mass at
entries whose density is less than \(\epsilon g\) is at most
\(\epsilon D\int g\).
These statements apply within conditional worlds and may be
summed with the original world weights.

Consequently further inverse-subpower restrictions, and refinements
by subpower lists, preserve typical logarithmic densities.
Upper density bounds may be integrated over arbitrary base regions
and summed over discrete entries after the exceptional observations
have been removed.
\end{lemma}

\begin{proof}
Restriction gives \(0\le g'\le g\), and integrating the defining
inequality on \(\{g'<\epsilon g\}\) proves
\Cref{a01:restriction-integral}.  For a refinement with densities
\(g'_c\), sum the inequalities \(g'_c<\epsilon g\) over at most \(D\)
labels at each old observation and then integrate.
Take \(\epsilon=N_0^{-\tau}\), and absorb a subpower \(D\) or a
subpower normalization loss into a smaller positive power.

For clarity about exceptional sets, write the pushed-forward
measure as a good part plus an omitted part, with densities
\(g_{\mathrm{good}},g_{\mathrm{bad}}\).
On the new states where \(g_{\mathrm{bad}}>g_{\mathrm{good}}\),
their total mass is at most \(2\int g_{\mathrm{bad}}\).
On the other states the total density is at most twice the good
density.  This transfers a ceiling proved for the restricted good
measure to a typical ceiling for the full resulting measure.
It does not justify summing an unremoved exceptional portion
against small conditional probabilities.
\end{proof}

A datum is \emph{known by \(p_s\)} if, after permitted pruning, it
belongs to a deterministic list of size at most \(K_0\) depending
only on \(p_s\).  List knowledge, rather than a choice of one
representative, is the convention throughout.
The exact base coordinates in an observation vary continuously.
To choose labels or density ranges while preserving temperedness, we
approximate their joint law on a sufficiently fine finite base mesh.
The next lemma controls the error uniformly over subsequent choices.

\begin{lemma}[Joint base flattening]\label{a01:flattening}
Fix tempered discrete labels and predicates of polynomial total
count.  Given a fixed \(T>0\), there is a fixed finite \(D\) such that
uniformly redistributing the base \((t,y)\) inside its mesh cells
of width \(\rho\sim N_0^{-D}\), while retaining the discrete entries,
changes the joint measure in total variation by \(O(N_0^{-T})\).
Thus a posterior list selector for these entries can be replaced,
with that error in success, by a cell-constant tempered list table.
\end{lemma}

\begin{proof}
Use coordinates \((t,y,V)\) and the remaining trajectory
coefficients, extending all densities by zero outside their
domains.  Holding the other variables fixed, each polynomial
condition has a bounded number of changes on a base-coordinate
axis, unless it is identically zero there; the latter condition
has no changes.  There are polynomially many conditions.
The weighted density and the discrete labels therefore change
across only polynomially many boundaries on each such fiber.
Their heights and the ranges of the remaining variables are
polynomially bounded.  Integration first on the fiber and then in
the other variables shows that a base translation of size at most
\(\rho\) changes the joint measure by at most
\(N_0^C\rho\), with \(C\) fixed by the tempered bounds.
Sum over the finitely many base axes and over internal index copies.

Within-cell averaging has the same bound, up to a dimensional
constant: write its \(L^1\) difference from the original density
as the integral of pairwise differences of density values in the
same cell, then compare these with translations of length at most
a constant times \(\rho\).  Choosing \(D>C+T+1\) proves the assertion.
All labels are included in this argument jointly.

A measurable selector has the same success up to total variation
error on the flattened law.  On that law the posterior of the
discrete entries is constant inside each cell with its conditioning
bins fixed; a constant maximizing list can be chosen there.
There are polynomially many such cells and entries.  The resulting
lookup table is tempered, and its success transfers back by the
same total variation estimate.
\end{proof}

When two grids overlap with at most \(K_0\) members at fixed exact
base and other conditioning data, their common refinement has the
same typical exponent by \Cref{a01:density-operations}.
For a one-way upper bound it suffices to sum the good measure over
the corresponding covering entries.

\begin{lemma}[Homogenization]\label{a01:homogenization}
After a tempered further selection of inverse-subpower mass and a
subsequence, an exact problem has a homogeneous profile \(f(r)\)
on every prescribed compact depth interval.  It is nondecreasing
and \(1\)-Lipschitz.  Jointly, one may homogenize the observation
\((p_r,[V]_{a_u})\) on compact nonnegative \((r,u)\)-ranges, obtaining
a function \(f(r,u)\) that is nondecreasing in each variable and
has Lipschitz constants \(1\) in \(r\) and \(j\) in \(u\).
Moreover \(f(r,0)=f(r)\).
\end{lemma}

\begin{proof}
Choose a countable dense depth grid containing all required
endpoints.  At stage \(N_0\), test a finite initial part of this
grid, increasing its size sufficiently slowly.
By \Cref{a01:flattening}, use a common fixed polynomial base mesh
fine enough for all the observations and for a power-small
total variation error.  The tested hidden and velocity grids
are nested and lie in fixed compact depth ranges, so their joint
entries are determined by the finest tested bins.  Their total
number, including base cells, is polynomial.

Polynomial bounds give a polynomial ceiling for each flattened
density.  Densities below \(N_0^{-C}\), for a sufficiently large
fixed \(C\), have power-small total mass, by summing over the
polynomial ranges and entries.  Quantize the remaining logarithmic
densities in intervals of width \(\epsilon_{N_0}\to0\).
If \(m_{N_0}\) is the number of tested observations, choose the
rates so slowly that the number of density-type lists is
\[
 \bigl(O(\epsilon_{N_0}^{-1})\bigr)^{m_{N_0}}=N_0^{o(1)}.
\]
One list carries inverse-subpower mass.  Restrict to it.
This restriction is a cell table testing the joint discrete
observations and is tempered.  Its upper bounds persist, and
\Cref{a01:density-operations} shows that its densities lose a
fixed power relative to the preselection values only on
power-small mass.

The total variation estimate transfers these conclusions to
the exact base densities.  Indeed, where two densities differ
by a large multiplicative factor, their \(L^1\) difference
controls the mass under the larger density there.
Pass to a subsequence on which the type exponents converge at
every fixed tested depth.

Refinement from depth \(r\) to \(r'\ge r\) has at most
\(K_0N_0^{r'-r}\) hidden-bin choices, and refinement from velocity
depth \(u\) to \(u'\ge u\) has at most
\(K_0N_0^{j(u'-u)}\) choices.  Monotonicity of densities under
refinement and \Cref{a01:density-operations} imply
\begin{align*}
 0&\le f(r',u)-f(r,u)\le r'-r,\\
 0&\le f(r,u')-f(r,u)\le j(u'-u).
\end{align*}
These inequalities extend the profiles uniquely to all
intermediate depths; nesting squeezes their typical densities
between nearby tested ones.  Since \(|V|\le K_0\), appending
the unit velocity bin costs only \(K_0\) choices, proving
\(f(r,0)=f(r)\).
\end{proof}

Other fixed logarithmic quantities with polynomial upper and
reciprocal bounds may be homogenized by the same argument.
In subsequent cap estimates, fixed tolerances on fixed finite
grids are sent to zero sufficiently slowly, and the grids
increase sufficiently slowly, that the exceptional mass is
negligible compared with the particular inverse-subpower
selection in use.  This convention concerns the present exact
\(N_0\)-sequence.

The class used in the contradiction imposes two further bounds.
The first limits how rapidly the typical density can decrease
as the hidden observation is refined toward the terminal depth.
The second limits velocity concentration within a terminal
observation.

\begin{definition}[The cap class]\label{a01:cap-class}
Fix \(S,\eta>0\).  A homogeneous exact problem belongs to \(C(d)\),
\(d<1\), if
\begin{equation}
 f(L)-f(r)\le d(L-r)\qquad(0\le r\le L),
 \label{a01:endpoint-cap}
\end{equation}
and, for \(0<u\le\eta L\), the joint density of
\((p_L,[V]_{a_u})\) has typical upper bound
\begin{equation}
 n(L)N_0^{-Su+o(1)}.
 \label{a01:angular-cap}
\end{equation}
The \(o(1)\) is interpreted by the fixed-tolerance convention above.
The parameters \(S,\eta\) are held fixed when the class is used.
\end{definition}

\begin{lemma}[Angular caps under coarsening]\label{a01:angular-coarsening}
Suppose \Cref{a01:angular-cap} holds. For \(0\le b<L\) and
\(0<u\le\eta L\), the conditional probability
\[
 \Pp\{[V]_{a_u}=v\mid p_b\},
 \qquad v=[V]_{a_u},
\]
under the selected incidence law has typical upper bound
\(N_0^{-Su+o(1)}\), evaluated at the sampled angular entry.
Normalizing the total incidence mass does not change this
conditional distribution. Good restricted measures with subpower
slack can be prepared simultaneously on slowly increasing finite
depth grids, losing negligibly relative to a specified
inverse-subpower selection.
\end{lemma}

\begin{proof}
Fix the exponents and a tolerance.  Remove the fine entries where
the joint angular upper bound or the pure \(p_L\) lower bound fails.
For each remaining fine entry and a specified angular bin, its
restricted numerator is bounded by
\(N_0^{-Su+O(\tau)}\) times the old pure density.
Sum these numerators over the fine hidden bins inside a coarse bin
before dividing by the corresponding sum of old pure densities.
This proves the coarse bound for the good submeasure relative to
the old denominator.  The omitted-mass comparison in
\Cref{a01:density-operations} gives the typical assertion, and
\Cref{a01:restriction-integral} handles a later inverse-subpower
restriction of the denominator.

For finitely many tests their power-small exceptions can be
summed.  Choose the tolerance and number of tests slowly to
obtain the final assertion.  These exclusions may be purely
analytical; if an output selection needs them,
\Cref{a01:flattening} implements the corresponding cell tests.
In particular this argument does not assert a bound on the
maximum angular probability over all bins of an unpruned law.
\end{proof}

\subsection{Time charts and prior-mass budgets}

We now describe the local fits whose time cost defines the horizon.
A chart controls each assigned trajectory throughout its time
interval. Its incidence selection specifies where the lower
derivative bound is required and supplies the mass entering the
true floor. This selection need not occupy the whole interval.

\begin{definition}[Charts, true packets, and known atlases]
\label{def:true-chart}
At depth \(r\in(0,L]\), a chart system has a common dyadic
time partition into intervals \(I\) of length
\(q=N_0^{-h+o(1)}\), \(h\ge0\).
Within each world and each interval, it assigns disjoint
trajectory sets \(A_c\) to chart labels \(c\), and selects
incidences from the given event on these assignments.
The total selected mass is at least \(K_0^{-1}\).

A chart has an affine spatial center \(y_c(t)\) and an
invertible matrix \(D_c\), with
\begin{equation}
 \|D_c\|\le qK_0,\qquad J_c=|\det D_c|,\qquad
 \sup_{t\in I}|D_c^{-1}(y(t)-y_c(t))|\le K_0
       \quad(l\in A_c).
 \label{a01:chart-space}
\end{equation}
It has a test of the appropriate model class satisfying
\begin{align}
 \sup_{t\in I}|P(t,y(t),w(t))|&\le K_0a_r,
 &
 \sup_{t\in I}|\partial_wP(t,y(t),w(t))|&\le K_0B,
 \label{a01:chart-test}\\
 |\partial_wP|&\ge B/K_0
       \quad\hbox{on its selected incidences},
 &
 |\partial_{ww}P|&\le K_0B^2/a_r.
 \notag
\end{align}
All chart data, assignments, selections, inverse coordinate
changes, and reciprocal widths are tempered.

Write
\[
 \nu_c=\nu_{\gamma(c)}(A_c),\qquad
 m_c=\mu_{\gamma(c)}(\hbox{selected incidences assigned to }c).
\]
A \emph{packet} is an individual chart together with its label,
time interval, geometric and test data, full trajectory assignment
\(A_c\), and current incidence selection. It carries the assigned
trajectory measure \(\nu_{\gamma(c)}|_{A_c}\), of mass \(\nu_c\),
and the unnormalized selected incidence measure, of mass \(m_c\).
An incidence-only restriction shrinks the latter without
conditioning the assigned trajectory measure on success.
A packet is \emph{true} when
\begin{equation}
 \frac{m_c}{q}\ge K_0^{-1}n(r)J_c.
 \label{a01:true-floor}
\end{equation}
A system is \emph{true} if every used packet is true.
A \emph{known atlas} at depth \(r\) is a chart system whose
label is known by \(p_L\). It has all the assignment, coverage,
geometric, test, and temperedness properties above, but need not
satisfy \Cref{a01:true-floor}.
We also allow known atlases at \(r=0\), with \(a_0=1\),
when the same displayed bounds hold.
\end{definition}

Since \(y-y_c\) is affine on \(I\), \Cref{a01:chart-space}
also bounds \(qD_c^{-1}(V-y_c')\) by \(K_0\).
Disjointness of assignments in each interval gives the basic
budget
\begin{equation}
 m_c\le q\nu_c,\qquad
 \sum_c\omega_{\gamma(c)}q\nu_c\le1.
 \label{a01:assignment-budget}
\end{equation}
The sum includes all worlds, intervals, and labels of one system.

The following elementary estimate converts the small-value and
lower-derivative tests into a bound on the hidden-coordinate bins.

\begin{lemma}[Quadratic sublevel sets]\label{lem:quadratic-sublevel}
Let \(Q\) be a real polynomial of degree at most two, and let
\(\varepsilon,b,\delta>0\). The set
\[
 \{w\in\R:|Q(w)|\le\varepsilon,
                  \ |Q'(w)|\ge b\}
\]
has total length at most \(4\varepsilon/b\) and meets at most
\(C(1+\varepsilon/(b\delta))\) bins of any grid of width
\(\delta\), for an absolute constant \(C\).
\end{lemma}

\begin{proof}
The condition \(|Q'|\ge b\) is the union of at most two
intervals on each of which \(Q\) is monotone. On either
interval, the inverse derivative bound gives length at most
\(2\varepsilon/b\) for the indicated sublevel set.
Each resulting interval meets at most two more grid bins than
its length divided by \(\delta\).
\end{proof}

The defining true floor compares incidence mass with observation
density and spatial volume. The next lemma combines that floor
with the assignment budget to obtain comparability with the prior
mass \(\nu_c\), and also confines the chart label to a subpower
list determined by the observation.

\begin{lemma}[Saturation and label knowledge]\label{a01:saturation}
A true system may be restricted, retaining inverse-subpower
total mass and enlarging subpower slacks, so that every used
chart satisfies
\begin{equation}
 \frac{m_c}{q}\sim_{\log}\nu_c\sim_{\log}n(r)J_c,
 \label{a01:saturated-floor}
\end{equation}
and its label is known by \(p_r\).
The true floor persists on the resulting selected incidences.
\end{lemma}

\begin{proof}
For a large subpower \(H\), deleting charts with
\(q\nu_c>Hm_c\) loses at most
\[
 \sum_{q\nu_c>Hm_c}\omega_{\gamma(c)}m_c
 \le H^{-1}\sum_c\omega_{\gamma(c)}q\nu_c\le H^{-1}.
\]
Choose \(H^{-1}\) negligible relative to the system mass.

Analytically remove the exceptional \(p_r\)-entries above the
ceiling \(K_0n(r)\), choosing the slack so the removed mass is
negligible relative to that same mass.
At a fixed base in chart \(c\), the conditions
\(|P|\le K_0a_r\) and \(|P_w|\ge B/K_0\) admit only \(K_0\)
hidden bins of width \(a_r/B\).
Apply \Cref{lem:quadratic-sublevel} with
\(\varepsilon=K_0a_r\), \(b=B/K_0\), and
\(\delta=a_r/B\), and enlarge the subpower factor.
The chart base volume is at most \(K_0qJ_c\).
Thus the good portion of \(m_c\) is at most \(K_0n(r)qJ_c\).
Discard whole charts whose full mass is more than twice this
bound.  Their total mass is controlled by twice the globally
omitted mass.  Together with the first deletion and the true
floor, this proves \Cref{a01:saturated-floor} before a final
enlargement of the slack.

To obtain lists, let a chart cross an observation if its time,
spatial, value, and event-derivative tests are compatible with
that observation.  If \(T(p_r)\) is the number of crossings,
the same volume calculation, now integrating the capped
original measure rather than only assigned incidences, gives
\begin{align*}
 \int T(p_r)\,d\mu_{\mathrm{cap}}
 &\le K_0\sum_c\omega_{\gamma(c)}n(r)qJ_c\\
 &\le K_0\sum_c\omega_{\gamma(c)}m_c.
\end{align*}
Markov's inequality, with a sufficiently large subpower
threshold, leaves subpower lists and loses negligible system
mass.  The joint flattening lemma implements these lists
as tempered tables if necessary.
Finally discard charts that lost more than half their
selected mass in these operations.  Their total loss is
at most twice the removed incidence mass, and
\Cref{a01:true-floor} survives with enlarged slack.
\end{proof}

Later restrictions can reduce the selected mass of an individual
chart. The next lemma gives two ways to continue a count: restore
a floor for the current entries by pruning, or charge earlier
witnesses using a cap for their own frozen measure.

\begin{lemma}[Restoring current floors and charging witnesses]
\label{a01:witness-budgets}
Fix a chart layer with the assignment budget
\Cref{a01:assignment-budget}.  Suppose a later selected measure
has inverse-subpower mass in the same normalization.
For each chart let at most \(D\le K_0\) additional entries be
used, and let \(\widetilde m_{c,e}\) be their current masses.
For any \(0<\epsilon<1\), deleting entries with
\begin{equation}
 \widetilde m_{c,e}<\frac{\epsilon q\nu_c}{D}
 \label{a01:current-floor}
\end{equation}
loses at most \(\epsilon\) total weighted mass.
In particular, \(\epsilon\) may be an inverse-subpower
negligible relative to the current total mass, giving
current floors against the full assignment budget.
For finitely many prescribed layers, allocate positive tolerances
\(\epsilon_i\) with \(\sum_i\epsilon_i\le\epsilon\). Their
corresponding floors can be achieved on one final retained measure,
with total weighted loss at most \(\epsilon\).

If time is further partitioned, the corresponding threshold
uses \(|I'|\nu_c\) for a subinterval \(I'\subset I\);
summing these budgets over the subintervals gives \(q\nu_c\).

For a collection of disjoint witness submeasures of a fixed
measure \(\lambda\), all supported in a domain \(U\), the total
witness mass is at most \(\lambda(U)\). With multiplicity at
most \(D\), the bound is \(D\lambda(U)\). Any cap used in this
count must hold for that same fixed measure \(\lambda\).
No floor on an arbitrary later restriction is thereby asserted.
\end{lemma}

\begin{proof}
Sum \Cref{a01:current-floor} over at most \(D\) entries per
chart, and then use \Cref{a01:assignment-budget}.
For finitely many layers, choose positive tolerances
\(\epsilon_i\) with \(\sum_i\epsilon_i\le\epsilon\), and use
\(\epsilon_i\) in the threshold for layer \(i\). Keep its assignment
budgets and entry counts fixed. Repeatedly delete all remaining
incidences of any entry whose current mass is positive and below its
fixed threshold,
reconsidering every layer after each deletion. An entry becomes empty
permanently when deleted, so it is charged at most once. The assignment
budget bounds the weighted charges from layer \(i\) by
\(\epsilon_i\), including deletions caused by earlier losses in other
layers. There are finitely many entries, so the process terminates
with every surviving entry meeting its threshold in the final measure
and total loss at most \(\epsilon\). The removed entries form a
polynomially sized table of original predicates, preserving temperedness.
For a slowly increasing collection satisfying the standing tempered
bounds, choose its size and tolerance allocation slowly enough that
the floor losses remain subpower.
The time-subinterval assertion follows from
\(\sum_{I'\subset I}|I'|\nu_c=q\nu_c\).

For the last assertion, if the witnesses of the counted
labels are disjoint submeasures of a frozen measure
\(\lambda\) and all lie in a domain \(U\), the sum of their
masses is at most \(\lambda(U)\).  A multiplicity bound \(D\)
replaces this by \(D\lambda(U)\).
The cap and the witnesses must therefore refer to the same
measure \(\lambda\), after every restriction required for that
cap has been imposed. A floor for these frozen witnesses gives no
lower bound for a later restriction; the first part of the lemma
provides a floor for that current measure.
\end{proof}

\begin{lemma}[Fullness of the largest spatial scale]
\label{a01:largest-axis}
In a cap-class problem, a chart system known by \(p_L\),
whether true or not, has
\[
 \|D_c\|\sim_{\log}q
\]
on an inverse-subpower selection, after homogenizing the
ratio if necessary.
\end{lemma}

\begin{proof}
The upper bound is part of \Cref{a01:chart-space}.
If the ratio were at most \(N_0^{-\alpha}\) for a fixed
\(\alpha>0\) on inverse-subpower mass, the endpoint bounds
in \Cref{a01:chart-space} would confine \(V\), for each
chart, to a ball of radius \(K_0N_0^{-\alpha}\) about \(y_c'\).
Choose a fixed \(u>0\) with \(u<\alpha/2\) and \(u\le\eta L\).
Such a ball meets only \(K_0\) velocity bins of width \(a_u\).
Since chart labels are known by \(p_L\), at each retained
terminal observation only \(K_0\) such velocity bins are
needed.  On the good restricted measure each has numerator
at most \(N_0^{-Su+o(1)}\) times the old pure denominator,
by \Cref{a01:angular-cap}.  Summing these bounds and then
integrating shows that the alleged mass is power-small.
The exceptional entries were removed before this sum.
This contradiction excludes every fixed-power deficiency.
\end{proof}

\subsection{Normalization, composition, and horizon}

To iterate the construction, we regard each chart as a new world
with its conditional trajectory prior and its own unit time
interval. The density transformation must be tracked within that
world before child charts can be lifted back to the original
problem. The following lemma records both operations and the
label-knowledge condition needed for intermediate true floors.

\begin{lemma}[Chart normalization and lifting]\label{lem:normalization}
Let a homogeneous exact problem of length \(L\) have a
terminal-known atlas at depth \(0\le r<L\).
After discarding extremely light charts and homogenizing
their Jacobians and prior masses, it defines a new
homogeneous exact problem of length \(L'=L-r\).
For its typical densities,
\begin{align}
 n'(L')&\sim_{\log}
       n(L)\frac{J_c}{\nu_c},
 \label{a01:terminal-normalization}\\
 n'(s)&\le N_0^{o(1)}
       n(r+s)\frac{J_c}{\nu_c}
       \qquad(0\le s<L').
 \label{a01:intermediate-normalization}
\end{align}
If the parent label is already known by \(p_r\), there is
equality up to subpower factors in
\Cref{a01:intermediate-normalization}.
The class \(C(d)\), with the same \(S,\eta\), is preserved.

A true terminal system in the new problem lifts to a true
terminal system in the old problem.  Intermediate true
systems also lift through a parent with the stated
intermediate density equality, in particular through a
prepared true parent.
\end{lemma}

\begin{proof}
Put
\[
 Z=\sum_c\omega_{\gamma(c)}q\nu_c,\qquad
 \omega'_c=\frac{\omega_{\gamma(c)}q\nu_c}{Z}.
\]
If the atlas success is \(M\ge K_0^{-1}\), then
\(M\le Z\le1\).  Use \(c=(\gamma,I,A_c)\) as a new world,
with trajectory law \(\nu_\gamma(\,\cdot\,|A_c)\) and
independent uniform time
\[
 t=t_c+q\tau,\qquad
 y=y_c(t)+D_cz,\qquad 0\le\tau\le1.
\]
The new selected measure is the pushforward of the old
assigned selected measure divided by \(q\nu_c\).
Its weighted total is the old atlas selected mass divided
by \(Z\).  In particular it is still a restriction of the
new product priors and has inverse-subpower mass.

In version \(1\), if the parent test is \(w-F(t,y)\), set
\[
 w'=\frac{w-F(t,y)}{a_r}.
\]
This is quadratic along trajectories and bounded by \(K_0\).
In version \(2\), retain \(w\) and use the native test
\(P_*'=P/a_r\), with derivative scale \(B'=B/a_r\).
The parent test bounds give
\[
 \sup|P_*'|\le K_0,\quad
 \sup|\partial_wP_*'|\le K_0B',\quad
 |\partial_wP_*'|\ge B'/K_0,\quad
 |\partial_{ww}P_*'|\le K_0(B')^2
\]
in the required places.  All base and velocity bounds
follow from \Cref{a01:chart-space}.
At \(r=0\), if the existing native test already has these
bounds after the spatial normalization, it may instead be
retained unchanged; this is useful for full-time
spatial atlases.

These coordinate changes preserve the permitted models.
Their maps on coefficient space have constant Jacobian:
the base changes are affine, and in version \(1\) the
subtraction of the quadratic \(F(t,y)\) is triangular
with respect to the \(w\)-coefficients.
They and their inverses have polynomial control.
Charts whose conditional prior is below a sufficiently
small reciprocal polynomial have negligible total
success, by polynomial chart count.  On the remaining
charts conditioning preserves temperedness.

At fixed old base and chart label, new hidden-bin width
at depth \(s\) is the old width \(a_{r+s}/B\), with a
chart- and base-dependent offset in version \(1\).
The two grids have bounded overlap.
For transported bins, the exact within-world density
transformation is
\begin{equation}
 \rho'_c
 =\frac{qJ_c}{q\nu_c}\rho_c
 =\frac{J_c}{\nu_c}\rho_c.
 \label{a01:exact-density-transform}
\end{equation}
The numerator \(qJ_c\) is the base Jacobian; the denominator
\(q\nu_c\) normalizes the new product prior.
Neither the world weight nor \(Z\) belongs in this
within-world factor.

Refining \(p_L\) by the known chart label does not change
its typical density exponent.  Apply
\Cref{a01:exact-density-transform}, the bounded bin overlap,
and \Cref{a01:density-operations} to obtain
\Cref{a01:terminal-normalization}.  Global exceptional
masses are multiplied by \(Z^{-1}\le K_0\), so their
power-small character is preserved.
At an earlier depth, refinement by an arbitrary label
can only decrease the unnormalized density, giving
\Cref{a01:intermediate-normalization}.  Knowledge by
\(p_r\), and hence by finer hidden observations, gives
the reverse typical inequality.
For a saturated true parent the formulas simplify to
\[
 n'(s)\sim_{\log}\frac{n(r+s)}{n(r)},\qquad
 n'(0)\sim_{\log}1.
\]

For the angular bound, a new velocity bin of width
\(N_0^{-u}\) maps by
\(V=y_c'+D_cV'/q\) into an old velocity ball of radius
\(K_0N_0^{-u}\).  At each exact chart and base it requires
only \(K_0\) old terminal/angular bins.
Sum the good old joint numerators before applying the
density transformation.  The old exceptional mass and
the new states where it dominates are handled as in
\Cref{a01:density-operations}.  This gives
\[
 \rho'_{\mathrm{joint}}\le
 N_0^{o(1)}n(L)\frac{J_c}{\nu_c}N_0^{-Su}
 \sim_{\log} n'(L')N_0^{-Su}
\]
for \(u\le\eta L'\).
Finally,
\[
 \log_{N_0}\frac{n'(s)}{n'(L')}
 \le f(L)-f(r+s)+o(1)
 \le d(L'-s)+o(1),
\]
which proves the endpoint cap.

For lifting, let a child have relative time length
\(q_d\), spatial Jacobian \(J_d\), and selected mass
\(m'_d\) in the new world.  Its old quantities are
\begin{equation}
 q_{cd}=qq_d,\qquad J_{cd}=J_cJ_d,\qquad
 m_{cd}=q\nu_c\,m'_d,\qquad
 \frac{m_{cd}}{q_{cd}}=
       \nu_c\frac{m'_d}{q_d}.
 \label{a01:composition-masses}
\end{equation}
At the terminal depth, a child true floor and
\Cref{a01:terminal-normalization} give the old floor
\(K_0^{-1}n(L)J_cJ_d\).
At an intermediate depth the same conclusion uses the
intermediate equality; an upper bound alone would not
suffice.

The tests lift by multiplication by \(a_r\).
In version \(1\), a child test \(w'-G(\tau,z)\) becomes
\[
 w-F(t,y)-a_rG(\tau,z),
\]
which is still of the special quadratic form.
In version \(2\), \(a_rQ(\tau,z,w)\) remains quadratic,
has old derivative scale \(B\), and satisfies
\[
 |\partial_{ww}(a_rQ)|
 \le K_0a_r\,\frac{(B/a_r)^2}{a_s}
 =K_0\frac{B^2}{a_{r+s}}.
\]
The small-value and first-derivative bounds transform
in the same way.  Spatial matrices and times compose,
whole-block fitting is preserved, and assignments may
record the entire path with nested time partitions.
All operations retain tempered bounds.
\end{proof}

In version \(2\), largest-axis fullness gives a geometric
interpretation to
\begin{equation}
 \log_{N_0}\frac{q^2}{J_c},
 \label{a01:narrowness}
\end{equation}
which we call the narrowness in the specified depth units:
it is the additional thinness of the smaller spatial axis,
up to subpower factors.
Under composition,
\[
 \frac{q_{cd}^2}{J_{cd}}
 =\frac{q^2}{J_c}\frac{q_d^2}{J_d}.
\]
Thus narrowness exponents add exactly, without any
assumption that the spatial axes are parallel.
Statements about a common value of this exponent use a
homogeneous subfamily.

\begin{definition}[Minimum horizon]\label{def:horizon}
For a homogeneous exact problem \(E\) of length \(L\), let
\(H(E)\) be the infimum of the exponents \(h\ge0\) for which
a true system at terminal depth \(L\) exists.
The infimum includes systems on arbitrary subsequences
and further permitted incidence selections, in the
coordinates of \(E\).  Such selections do not change
the terminal density exponent.
\end{definition}

The rest of the section proves two complementary bounds.
Every homogeneous exact problem satisfies \(H(E)\le L\):
terminal fits can be built on intervals of length comparable to
\(a_L\).
A hypothetical Kakeya counterexample, however, yields a cap-class
problem with \(H(E)>0\). Both arguments use the weighted plank
estimate below.

\subsection{The weighted plank input and the initial horizon bound}

We first prove that true terminal systems always exist with horizon
at most \(L\). At time scale \(a_L\), the native quadratic signal
already supplies the polynomial test. A weighted plank estimate
will supply spatial boxes whose incidence mass meets the true floor.

\begin{lemma}[Weighted full-time plank estimate]
\label{lem:weighted-planks}
Fix a bounded affine-line chart in \(\R^2\), and a fixed
permitted enlargement of mesh cells.
For every \(\eps>0\), there are \(\eta_\eps>0\) and
\(C_\eps<\infty\) with the following property.
Let
\[
 M_i(t)=b_i+tu_i,\qquad 0\le t\le1,
\]
be finitely many indexed traces, with \(b_i,u_i\) in the
fixed bounded chart, and arbitrary positive weights
\(\omega_i\).  Repeated traces are allowed.
Let \(D\ge1\), \(\rho=D^{-1}\), and suppose that every
full-time plank
\[
 \mathcal P=\left\{(t,x):
 \begin{array}{l}
 0\le t\le1,\\
 |\langle x-c-tv,e_1\rangle|\le a\rho,\\
 |\langle x-c-tv,e_2\rangle|\le b\rho
 \end{array}\right\},
 \qquad 1\le a\le b\le D,
\]
where \((e_1,e_2)\) is an arbitrary orthonormal spatial
frame and \(c,v\in\R^2\), satisfies
\begin{equation}
 \sum_{i:\,M_i\subset\mathcal P}\omega_i\le Aab.
 \label{a01:weighted-plank-cap}
\end{equation}
Fixed dilations of these plank tests are permitted.
Mark time bins \(J_i\), each of length \(\rho\), with
\(\#J_i\ge D^{1-\eta_\eps}\).
If \(E_{\mathrm{vis}}\) is the union of the spatial--time
mesh cells visited at their bin centers, or fixed
enlargements of those cells, then
\begin{equation}
 \sum_i\omega_i\#J_i
 \le C_\eps A D^\eps\,\#E_{\mathrm{vis}}.
 \label{a01:weighted-plank-conclusion}
\end{equation}
\end{lemma}

\begin{proof}
This is the indexed weighted formulation of
\citet[Lemma~2.3]{ThreeDimensional}.
The plank condition concerns containment of the entire
trace, not merely its marked portions.
The cited statement permits arbitrary positive weights
and repeated indices, and has precisely the long-mark
hypothesis above.
\end{proof}

We record explicitly two extensions of use below.
A continuous weighted law of bounded affine traces can be
rounded in coefficient space at a mesh much smaller than
\(\rho\).  Group traces with the same rounded coefficients
and the same finite mark set, and assign each group its
total prior mass.  This gives a finite indexed family for
\Cref{lem:weighted-planks}.  Fixed enlargements of cells and
planks recover the original traces, and do not alter the
estimates beyond constant factors.  These finite
approximations are used for the estimate; their mark
sets need not be output labels of a chart construction.

Also suppose the chart parameters are bounded by a
subpower \(K_0\), the mesh is a fixed power of \(N_0^{-1}\),
and the available labeled lengths are at least \(K_0^{-1}\).
A common spatial rescaling by a subpower places the traces
in a fixed bounded chart.  Refining meshes and enlarging
cells by subpower factors changes all mark and cell
counts by subpower factors.  The long-mark condition is
then satisfied for every fixed \(\eps>0\) at sufficiently
large \(N_0\).  Choose \(\eps\downarrow0\) only after these
fixed-\(\eps\) requirements, by a slow diagonal.
The factors \(C_\eps D^\eps\) are thereby subpower.
In physical widths \(u,v\), the resulting conclusion
bounds raw weighted average multiplicity by a subpower
times the supremum of
\(\rho^2\sum_{M_i\subset\mathcal P}\omega_i/(uv)\).
This is the form used when a plank is inferred from a
small visited-cell set.

\begin{proposition}[Initial horizon bound]\label{a01:horizon-bound}
Every homogeneous exact problem satisfies \(H(E)\le L\).
\end{proposition}

\begin{proof}
Choose a dyadic \(q\sim a_L\).
In each world and each \(q\)-interval \(I\), group
trajectories by \(y\) at the midpoint of \(I\), to width
\(q\), and by \(X\) there, to width \(a_L\).
Call the joint label \(G\).
Along the block, both \(y\) and \(X\) move by at most
\(K_0q\): for \(X\), use its quadratic coefficients and
the uniform bound in the exact problem.
At an observed incidence, the hidden bin determines \(X\)
to \(K_0a_L\).  In version \(2\), Taylor's formula in \(w\)
uses the native first-derivative and curvature bounds;
the curvature contribution is \(O(K_0a_L^2)\).
Thus \(G\) is known by \(p_L\) up to subpower lists.

Let \(b_G\) be the \(X\)-bin center.
Use the test \(w-b_G\) in version \(1\), or \(P_*-b_G\)
in version \(2\).
It has the required terminal small-value and derivative
bounds throughout the block and at selected incidences.
It remains to find spatial planks carrying the true mass
floor.

Translate by the spatial bin center and divide spatial
coordinates by \(Rq\), where \(R\) is a sufficiently
large subpower.  Normalize block time to \([0,1]\).
All normalized affine traces then lie in one fixed
bounded chart.  Put \(J_0=(Rq)^j\).
In these coordinates use the old trajectory measure
inside \(G\), without normalizing it to a probability;
only time is normalized to the block.
The typical exact-base density in this group is at least
\(n(L)J_0\) up to subpower factors: refine by the known
label \(G\), change coordinates, and sum over compatible
hidden bins.

Choose a common mesh \(\rho=D_m^{-1}\sim N_0^{-D}\),
with \(D_m\) dyadic.  Here \(D\) is fixed sufficiently
large for the tempered bounds and total variation
errors, before the loss parameters below are chosen.
Joint flattening transfers the preceding lower bound
to the cell averages, apart from power-small exceptional
mass, weighted over all worlds, groups, and blocks in
original time.  Coordinate changes and block
normalizations have only polynomial costs.

Fix \(0<\xi<1\).
In each group greedily select spatial planks with
arbitrary time-independent orthonormal spatial frames, affine centers, and
normalized half-widths
\(\rho\le r_i\lesssim1\), \(1\le i\le j\).
A plank may be selected if the remaining incidence mass
on trajectories wholly contained in it, with block time
normalized, is at least
\begin{equation}
 N_0^{-\xi}n(L)J_0\prod_{i=1}^j r_i.
 \label{a01:greedy-threshold}
\end{equation}
Assign all such available trajectories to that chart and
remove them from the group.  Fixed constants in the
width cutoffs are chosen to allow subsequent
enlargements.

This procedure has polynomially many steps, uniformly
for \(0<\xi<1\): the mass removed in any step is bounded
below by a reciprocal polynomial, and the total
within-world block mass is at most one.
Full containment is tested at endpoints in each spatial
axis, so assignments and successive exclusions are
tempered.  A nonempty plank can have its center and
speed controlled by choosing a contained bounded trace
as center and enlarging the widths by a fixed factor.
Its old spatial Jacobian is \(J_0\prod_i r_i\), and the
test remains the one fixed by \(G\).

Suppose the algorithm covers less than half the initial
selected mass along a subsequence.
In the residual remove cells that fail the original
group floor with tolerance \(\xi/4\), then remove
trajectory--block entries whose remaining normalized
labeled duration is less than the initial total selected
mass divided by a sufficiently large constant.
The first loss is power-small.  The second is bounded
by that duration threshold, because the trajectory
assignments and world--block weights have total prior
budget at most one.
There remains inverse-subpower mass, and every used
trajectory has labeled normalized duration at least
\(K_0^{-1}\).

Write \(m_G\) for the original normalized incidence mass
of a group.  Weighted averaging gives a group with
surviving normalized mass at least \(m_G/K_0>0\).
Every residual visited good cell has original mass at
least
\[
 N_0^{-\xi/4}n(L)J_0\rho^{j+1}.
\]
The number of such cells is consequently at most
\begin{equation}
 \frac{m_G}
 {N_0^{-\xi/4}n(L)J_0\rho^{j+1}}.
 \label{a01:old-cell-count}
\end{equation}
The mass used in this count is the original whole-cell
witness.  A residual cell need not retain that floor.

Mark all time bins with positive surviving labeled
duration on each used trace.  There are at least
\(D_m/K_0\) marks per trace, so
\Cref{lem:weighted-planks} applies for any fixed
\(\eps>0\), eventually.
Moving a point to the bin center costs a fixed cell
enlargement because normalized speeds are bounded.
Give each trace its old trajectory weight.

For \(j=1\), add an independent static coordinate
uniform on \([0,1]\); this increases the visited-cell
count by at most \(O(\rho^{-1})\) and preserves total
weighted mark count.  In both cases, using
\Cref{a01:old-cell-count},
\begin{equation}
 \frac{\sum_i\omega_i\#J_i}{\#E_{\mathrm{vis}}}
 \gtrsim
 K_0^{-1}N_0^{-\xi/4}n(L)J_0\rho^2,
 \label{a01:greedy-raw-multiplicity}
\end{equation}
since the numerator is at least \(m_GD_m/K_0\).

Taking the supremum plank constant in
\Cref{a01:weighted-plank-cap}, the weighted estimate
therefore supplies a full-time normalized plank of
physical spatial area \(uv\) with
\begin{equation}
 \frac{\hbox{used trajectory weight in the plank}}{uv}
 \gtrsim_\eps
 K_0^{-1}N_0^{-\xi/4}D_m^{-\eps}n(L)J_0.
 \label{a01:greedy-heavy-plank}
\end{equation}
The harmless fixed enlargement also contains the
unrounded traces.

In the padded case, project the rectangle onto the true
spatial coordinate, and denote its projected width by
\(W\).  The maximal vertical section length of a
rectangle of area comparable to \(uv\) is at most
\(Cuv/W\).
For example, writing its side half-widths as \(u,v\)
and its angle as \(\theta\), its projected half-width is
\(u|\cos\theta|+v|\sin\theta|\), while a vertical chord
has length at most
\[
 2\min\!\left(\frac{u}{|\sin\theta|},
             \frac{v}{|\cos\theta|}\right),
\]
with the evident interpretation at zero denominators.
Their product is at most a fixed multiple of \(uv\).
For each original trace the accepted fraction of the
static dummy coordinate is no larger than this
section bound, already at one time.
Thus projection turns \Cref{a01:greedy-heavy-plank} into
the same lower bound for original trace weight divided
by \(W\).

Every counted original residual trace has labeled
duration at least \(K_0^{-1}\).  In either dimension,
the resulting original incidence mass in the projected
or unpadded plank is therefore larger than the greedy
threshold, for large \(N_0\), if
\[
 D\eps<\xi/4.
\]
The two fixed losses then total less than \(\xi/2\);
all other losses are subpower.
This contradicts the stopping rule.

Consequently at least half the selected mass is covered.
Let \(\xi\downarrow0\) sufficiently slowly, after each
fixed-\(\xi,\eps\) estimate is valid.
The polynomial complexity bounds above were uniform
for \(\xi<1\), so the output remains tempered.
The threshold \Cref{a01:greedy-threshold} is now the
true floor with a subpower loss, and \(q\sim a_L\)
gives horizon exponent \(h=L\).
\end{proof}

\subsection{From an arbitrary Kakeya counterexample to positive horizon}

\begin{theorem}[Reduction to a bad weighted problem]
\label{thm:model-reduction}
If a set \(K\subset\R^4\), with no measurability or
compactness assumption, contains a unit segment in every
direction and satisfies \(\dim_{\mathrm H}K<4\), then
there are \(0<d<1\), \(S,\eta>0\), and a tempered
homogeneous exact problem \(E\) of version \(2\) in
\(C(d)\) such that \(H(E)>0\).
\end{theorem}

\begin{proof}
We give the selection from \(K\) first, then construct
the cap-class problem, and finally prove the positive
horizon bound.

\paragraph{A finite selection from outer measure.}
Fix a bounded open graph-slope chart of directions,
with slopes \(s\in\R^3\), writing spatial lines as
\[
 x(t)=b+ts,\qquad x=(y,w).
\]
A unit segment in one of these directions has a
nondegenerate time projection.  Every such segment
contains the graph over some interval with rational
endpoints, and its intercept belongs to a bounded box
with integer endpoints.  These are countably many
choices.  Countable outer subadditivity therefore
supplies one interval and one intercept box for which
the set \(S_0\) of eligible slopes has positive outer
Lebesgue measure, say \(a>0\).
After a fixed affine coordinate change, the common
time interval is \([0,1]\), and all these slopes and
intercepts lie in fixed bounded boxes.
This coordinate change preserves Hausdorff dimension.
No direction-to-witness map has been chosen.

Choose \(0<\sigma<1\) with
\(\dim_{\mathrm H}K<4-\sigma\).
There are covers by balls of arbitrarily small
maximum radius with
\[
 \sum_i r_i^{\,4-\sigma}\le1.
\]
Only balls meeting the fixed bounded region containing
the eligible graph segments are needed for the
following argument; their centers are bounded as well.
For \(k\ge1\), let \(U_k\) be the union of balls with
\(2^{-k-1}<r_i\le2^{-k}\).
There are at most \(C2^{k(4-\sigma)}\) such balls.

Choose a fixed \(c>0\) so that
\(\sum_{k\ge1}ck^{-2}<1\).
Every eligible witness has some \(k\) for which
\[
 |\{t\in[0,1]:b+ts\in U_k\}|\ge ck^{-2},
\]
because the sum of these time measures is at least one.
Define \(S_k\subset S_0\) by existence of a bounded
witness with this property.
The sets \(S_k\) need not be measurable, but
\(S_0\subset\bigcup_kS_k\).
Choose \(c'>0\) with
\(c'\sum_{k\ge1}k^{-2}<1\).
Outer subadditivity implies that for some \(k\)
\begin{equation}
 |S_k|^*\ge c'a k^{-2}.
 \label{a01:productive-slopes}
\end{equation}
Indeed the reverse strict inequality for every \(k\)
would imply \(|S_0|^*<a\).
The same scale thus has both the time-fraction
condition defining \(S_k\) and
\Cref{a01:productive-slopes}.
As the maximum cover radius tends to zero, every
possible such \(k\) tends to infinity.

Put \(\delta=2^{-k}\).
A maximal \(\delta\)-separated subset of the bounded
set \(S_k\) is finite and covers \(S_k\) by
\(\delta\)-balls.  Its cardinality \(M\) therefore
satisfies
\begin{equation}
 c k^{-2}\delta^{-3}\le M\le C\delta^{-3}.
 \label{a01:finite-slope-packing}
\end{equation}
Choose one successful witness for each of these
finitely many slopes.  This is the only witness
choice needed, and uses no measurable selection.

\paragraph{Smoothing and three density bounds.}
Let \(N_0=\delta^{-1}\), initially \(L=1\),
\(P_*=w\), and \(B=1\).
Use the uniform mixture of the selected lines and
perturb all three intercept and all three slope
coordinates by independent uniform boxes of radius
a fixed small multiple of \(\delta\).
Select incidences in fixed enlargements of the balls
in \(U_k\).
Each original covered time stays in those enlargements
under every allowed perturbation, so selected mass is
at least \(ck^{-2}\).

This is a tempered version-\(2\) problem.
The mixture may be represented by the internal indexed
copies of \Cref{a01:tempered}; its count and density
heights are polynomial, its supports are boxes, and
the finite union of enlarged balls is a bounded-degree
polynomial predicate of polynomial size.
The separated slope centers and
\Cref{a01:finite-slope-packing} give the full
three-slope density ceiling
\begin{equation}
 \frac{d\nu_s}{ds}\le Ck^2\le K_0.
 \label{a01:original-slope-cap}
\end{equation}
Conditional on the full slope, the \(y\)-intercept
density is at most \(C\delta^{-2}\): within each
mixture component this follows from independent
intercept smoothing, and conditional mixture weights
sum to one.

Homogenize jointly for \(0\le r\le1\), \(0\le u\le4\),
where \(u\) is the depth of the two-dimensional
\(V\)-bin.  Write
\[
 n(r,u)=N_0^{-f(r,u)},\qquad
 m(r,u)=N_0^r n(r,u),\qquad
 g(r,u)=r-f(r,u).
\]
The resulting limiting exponents obey
\begin{equation}
 g(1,0)\ge\sigma,\qquad
 g(0,0)\le0,\qquad
 g(r,u)\le r+2-2u.
 \label{a01:initial-density-inequalities}
\end{equation}

For the first inequality, at width \(\delta\) each
enlarged covering ball uses \(O(1)\) hidden \(w\)-bins
over a base \((t,y)\)-volume \(O(\delta^3)\).
The total observation-space measure of these states
is at most \(C\delta^{\sigma-1}\).
The selected mass is inverse-subpower and the typical
density is \(n(1,0)\), so
\(n(1,0)\ge K_0^{-1}\delta^{1-\sigma}\).
This gives \(g(1,0)\ge\sigma\).

For the third inequality, the \(V\)-marginal has bounded
region and density at most \(K_0\), by integrating
\Cref{a01:original-slope-cap} over the bounded remaining
slope.  A \(V\)-bin has probability at most
\(K_0N_0^{-2u}\).
At fixed \(t\), the conditional \(y(t)\)-density,
given the full slope, is at most \(C\delta^{-2}\).
Ignoring the hidden \(w\)-bin gives the pointwise
joint ceiling \(K_0N_0^{2-2u}\), hence the claimed
bound for \(g(r,u)\).

For the second inequality, the projected \(y\)-lines
have a full-time plank cap at every fine mesh
\(\rho=D_m^{-1}\):
\begin{equation}
 \nu\{l:y_l([0,1])\subset\mathcal P\}
 \le K_0ab\rho^2
 \label{a01:projected-full-plank-cap}
\end{equation}
for spatial radii \(a\rho,b\rho\).
Containment at \(t=0,1\) confines the velocity to the
corresponding rectangle of radii \(2a\rho,2b\rho\),
and its density is at most \(K_0\).

Suppose instead that \(n(0,0)\) grows by a fixed
positive power.  Choose \(\alpha>0\) smaller than that
power, and flatten at a sufficiently fine fixed
polynomial mesh \(\rho\sim N_0^{-D}\).
Retain typical cells.  Their projected base cell
masses have the lower bound
\(N_0^\alpha\rho^3\), after decreasing \(\alpha\)
if needed; the bounded \(w\)-range costs only
subpower many unit bins.
There are therefore at most
\(K_0N_0^{-\alpha}\rho^{-3}\) visited base cells.
Prune short labeled durations using the product
prior, leaving inverse-subpower total mass and
duration at least \(K_0^{-1}\) per used trace.
Write \(m_{\mathrm{ret}}\ge K_0^{-1}\) for the retained
incidence mass, and mark the visited time bins of each trace.
Each mark covers at most \(\rho\) units of time. In the finite
weighted approximation used for \Cref{lem:weighted-planks}, this gives
\[
 \sum_i\omega_i\#J_i\ge\frac{m_{\mathrm{ret}}}{\rho}.
\]
Dividing by the visited-cell bound, and enlarging the subpower
factor, yields the raw weighted average multiplicity
\[
 \frac{\sum_i\omega_i\#J_i}{\#E_{\mathrm{vis}}}
 \ge\frac{m_{\mathrm{ret}}/\rho}
              {K_0N_0^{-\alpha}\rho^{-3}}
 \ge K_0^{-1}N_0^\alpha\rho^2.
\]
On the other hand,
\Cref{lem:weighted-planks,a01:projected-full-plank-cap}
bound it by \(C_\eps K_0D_m^\eps\rho^2\).
Choose \(D\eps<\alpha/2\).
This is impossible for large \(N_0\), proving
\(g(0,0)\le0\).

\paragraph{Selecting an endpoint and an angular scale.}
The bounds on \(g\) force positive growth between hidden depths
zero and one and give an upper bound at large velocity depth.
We use both facts to select an endpoint with the one-sided
comparisons required by the cap class.
Choose small constants \(R,S>0\), for example with
\(R+4S<\sigma/2\) and \(R,S<1/10\), and maximize
\begin{equation}
 g(r,u)-Rr+Su
 \quad\hbox{on }[0,1]\times[0,4].
 \label{a01:endpoint-maximization}
\end{equation}
Continuity follows from homogenization.
At a maximizer \((r_0,u_0)\), we have
\begin{equation}
 r_0>0,\qquad u_0<4,\qquad
 g(r_0,u_0)\ge\sigma/2.
 \label{a01:chosen-endpoint}
\end{equation}
Indeed the value at \((1,0)\) is at least \(\sigma-R\).
At \(r=0\), monotonicity in \(u\) and
\(g(0,0)\le0\) bound the value by \(4S\).
At \(u=4\), the last inequality in
\Cref{a01:initial-density-inequalities} bounds it by
\(r-6-Rr+4S<0\).
Finally maximality gives
\(g(r_0,u_0)\ge\sigma-R-4S\).

Partition the \(y\)-intercept and \(V\) coordinates into
full-time boxes \(\mathcal B\) of side
\(a_{u_0}\), with dyadic rounding.
Make them worlds with conditional trajectory priors
and weights proportional to their original
probabilities \(\nu(\mathcal B)\).
Discard extremely light boxes and homogenize their
weights.
Subtract the central affine \(y\)-line \(y_{\mathcal B}(t)\) and divide
\(y\) by \(a_{u_0}\), retaining \(w\).
At fixed exact \((t,y)\), the old \(V\)-bin determines the
intercept bin up to boundedly many possibilities, since
\(y_0=y-tV\). Thus appending the crop label \(\mathcal B\)
to \((p_s,[V]_{a_{u_0}})\) is a bounded-list refinement at
every depth \(s\).

For clarity, let \(\rho_{\mathcal B}(s)\) be the density of
this refined observation in the old, unnormalized law.
Time is unchanged, the base change
\(y=y_{\mathcal B}(t)+a_{u_0}y_{\mathrm{new}}\) has
Jacobian \(a_{u_0}^2\), and conditioning the trajectory
prior divides by \(\nu(\mathcal B)\). For transported
bins the within-world density is therefore exactly
\[
 \rho_{\mathrm{new},\mathcal B}(s)
 =\frac{a_{u_0}^2}{\nu(\mathcal B)}\rho_{\mathcal B}(s).
\]
This is the specialization of
\Cref{a01:exact-density-transform} to a full-time crop.
The actual rounded grids have bounded overlap. Combining that
overlap and the bounded crop-label ambiguity with
\Cref{a01:density-operations} gives the typical equivalence,
for \(0\le s\le r_0\),
\begin{equation}
 n_{\mathrm{new}}(s)\sim_{\log}
 n(s,u_0)\frac{a_{u_0}^2}{\nu(\mathcal B)}.
 \label{a01:crop-density}
\end{equation}
A new angular bin of depth \(u\) corresponds, with the same
bounded-overlap convention, to the old joint angular refinement
at depth \(u_0+u\).

Maximality in \Cref{a01:endpoint-maximization} gives
\begin{align}
 f(r_0,u_0)-f(s,u_0)
   &\le(1-R)(r_0-s),&&0\le s\le r_0,
 \label{a01:crop-cap}\\
 f(r_0,u_0+u)-f(r_0,u_0)
   &\ge Su,&&0\le u\le4-u_0.
 \notag
\end{align}
Consequently the new exact problem, with length
\(L=r_0\), belongs to \(C(1-R)\) for the chosen \(S\)
and, for example,
\(\eta=(4-u_0)/(2r_0)>0\).

It remains to prove positive horizon. We return a true terminal
chart's mass floor to the original prior, then bound the prior
mass of all trajectories that can satisfy its geometric tests.
Comparing these bounds will force its time interval to be short.

\paragraph{The original-prior floor of a terminal chart.}
Consider any true terminal chart in this new problem,
including one obtained after a further permitted
selection or on a subsequence.
Let its time length be \(q\), its spatial Jacobian
in the cropped coordinates be \(J_{\mathrm{new}}\),
and set
\[
 J_{\mathrm{phys}}=a_{u_0}^2J_{\mathrm{new}}.
\]
Multiplying its true floor by \(\nu(\mathcal B)\)
and using \Cref{a01:crop-density}, its selected mass
in the original, unnormalized prior satisfies
\begin{equation}
 \frac{m_{\mathrm{old},c}}q
 \ge K_0^{-1}n(r_0,u_0)J_{\mathrm{phys}}
 =K_0^{-1}m(r_0,u_0)a_{r_0}J_{\mathrm{phys}}.
 \label{a01:original-chart-floor}
\end{equation}
This calculation cancels \(\nu(\mathcal B)\) even
when the crop has a very small polynomial mass.
Neither the crop prior nor the original prior is
renormalized to the analytically successful incidences.

Every trajectory contributing positive selected time
to this floor satisfies the whole-block spatial,
small-value, and upper-derivative tests in the original
coordinates, and satisfies \(|P_w|\ge K_0^{-1}\)
at some time.  The original prior mass of these
trajectories is at least \(m_{\mathrm{old},c}/q\).
Unused trajectories in the chart assignment are not
needed for this lower bound.

\paragraph{The slope-volume bound.}
Fix this chart.
Let \(\Sigma\subset\R^3\) consist of the bounded slopes
\(s\) for which some bounded intercept \(b\) satisfies
\[
 \sup_{t\in I}|D_{\mathrm{phys}}^{-1}
                (y_{b,s}(t)-y_c(t))|\le K_0,\qquad
 \sup_{t\in I}|P(t,b+ts)|\le K_0a_{r_0},
\]
and
\[
 \sup_{t\in I}|P_w(t,b+ts)|\le K_0,\qquad
 \sup_{t\in I}|P_w(t,b+ts)|\ge K_0^{-1}.
\]
Here \(P(t,b+ts)\) means that the three spatial
coordinates are split as \((y,w)\), and
\(D_{\mathrm{phys}}\) is the chart's original spatial
\(y\)-matrix.
Use fixed bounded slope and intercept boxes containing
the original smoothed prior support.

The set \(\Sigma\) is defined by these geometric tests alone.
We impose neither crop membership, finite witness identity, density
selection, nor selected-time conditions. Thus \(\Sigma\) contains
the slopes of all trajectories contributing to
\Cref{a01:original-chart-floor}.

We claim that
\begin{equation}
 |\Sigma|\le K_0\frac{a_{r_0}J_{\mathrm{phys}}}{q^3}.
 \label{a01:eligible-slope-volume}
\end{equation}
The original full-slope density ceiling will then convert this
geometric bound into an upper bound for the prior mass of an
eligible chart. To prove the claim, we select one eligible
intercept for each slope and compare the volume swept out by
these lines with the volume allowed by the chart tests.

The definition of \(\Sigma\) uses a fixed number of
real variables and fixed-degree polynomial conditions,
with quantification over \(b\) and \(t\).
All chart coefficients, bounds, and interval endpoints
are parameters.
Quantifier elimination, semialgebraic choice, and
\(C^1\) cell decomposition therefore provide an
eligible choice \(b=b(s)\) whose graph has uniformly
bounded formula complexity; its \(C^1\) full-dimensional
open pieces cover \(\Sigma\) up to a null set.
We use here the parameter-uniform forms of these
foundational results from real algebraic geometry;
see \citep{BasuPollackRoy}.
Their bounds depend on the formula format, not on
the numerical magnitudes of the chart coefficients.

For each \(t\), put \(F_t(s)=b(s)+ts\).
On a \(C^1\) piece,
\begin{equation}
 \det DF_t(s)=\det(Db(s)+t\,\mathrm{Id}),
 \label{a01:monic-jacobian}
\end{equation}
a monic cubic in \(t\).
There is also a uniform bound on the number of regular
preimages of any \(F_t\).
Indeed the fiber has fixed semialgebraic complexity,
uniformly in \(t\), the target point, and the chart
parameters, so has a uniformly bounded number of
connected components.  A preimage with nonzero
Jacobian is isolated in its fiber by local inversion
inside its open \(C^1\) piece, and hence is a separate
component.  This proves the bound without any bound
on \(Db\), and excludes a multiplicity growing with
\(N_0\).

\paragraph{Times with a large Jacobian and derivative.}
For a fixed slope \(s\), remove from \(I\) intervals
of radius \(cq\) around the real parts of the three
complex roots of the monic cubic in
\Cref{a01:monic-jacobian}.
The total removed length is at most \(6cq\).
Outside them each factor has modulus at least \(cq\),
so
\[
 |\det DF_t(s)|\ge c^3q^3.
\]
Also,
\(g_s(t)=P_w(t,b(s)+ts)\) is real affine and has
\(\sup_I|g_s|\ge K_0^{-1}\).
For small fixed \(\epsilon>0\),
\begin{equation}
 |\{t\in I:|g_s(t)|<\epsilon/K_0\}|
 \le4\epsilon q.
 \label{a01:affine-good-time}
\end{equation}
To verify this, let \(M=\sup_I|g_s|\).
If the variation of \(g_s\) on \(I\) is less than
\(M/2\), then \(|g_s|\ge M/2\) throughout.
Otherwise its slope has modulus at least \(M/(2q)\),
and the inverse image of
\((-\epsilon/K_0,\epsilon/K_0)\) has length at most
\(4\epsilon q\).

Choose \(c,\epsilon\) so that each bad set has
length at most \(q/4\).
Let \(\Sigma_t\) be the slopes in the open
\(C^1\) pieces passing both tests.
Then
\begin{equation}
 \int_I|\Sigma_t|\,dt\ge\frac q2|\Sigma|.
 \label{a01:good-time-overlap}
\end{equation}
The two good-time conditions therefore overlap
uniformly for each slope; no independence of them
is asserted.

The area formula and the uniform regular-fiber bound
give
\begin{align}
 q^4|\Sigma|
 &\lesssim q^3\int_I|\Sigma_t|\,dt
 \lesssim \int_I|F_t(\Sigma_t)|\,dt\\
 &\le K_0qJ_{\mathrm{phys}}a_{r_0}.
 \label{a01:selector-volume}
\end{align}
For the last inequality, the allowed \(y\)-region
has area at most \(K_0J_{\mathrm{phys}}\).
At fixed \(t,y\), the conditions
\[
 |P(t,y,w)|\le K_0a_{r_0},\qquad
 |P_w(t,y,w)|\ge\epsilon/K_0
\]
confine \(w\) to total length at most
\(K_0a_{r_0}\), after enlarging the subpower factor.
This is \Cref{lem:quadratic-sublevel}, applied at fixed
\((t,y)\).
Integrate this spatial-volume bound over \(I\). Dividing
\Cref{a01:selector-volume} by \(q^4\) proves
\Cref{a01:eligible-slope-volume}.

\paragraph{Comparison with the mass floor.}
Use the \emph{original} full-slope density
ceiling \Cref{a01:original-slope-cap}.
The original prior mass of all geometrically eligible
trajectories is at most \(K_0|\Sigma|\).
Combining this with
\Cref{a01:original-chart-floor,a01:eligible-slope-volume}
and cancelling \(a_{r_0}J_{\mathrm{phys}}>0\) yields
\[
 m(r_0,u_0)\le K_0q^{-3}.
\]
By \Cref{a01:chosen-endpoint},
\(m(r_0,u_0)\ge N_0^{\sigma/2}\).
Thus any true terminal system with
\(q=N_0^{-h+o(1)}\) has \(h\ge\sigma/6\).
The same estimate holds on every subsequence and
further allowed selection, so
\(H(E)\ge\sigma/6>0\).
This proves the theorem.
\end{proof}

We call a cap-class problem with \(H(E)>0\) bad.
The remaining sections exclude bad problems for every
fixed \(S,\eta>0\), first in version \(1\) and then in
version \(2\).  By \Cref{thm:model-reduction}, their
exclusion proves the theorem for arbitrary Kakeya sets.
